\begin{theorem}\label{fam:src:D.2.3}\label{fam:st:D.2.3}
The following descriptions give all compact irreducible curves of $X_c$.
In $W_c$ and $S_j$ they are exactly the images described in \ref{fam:src:D.2.2}.
In a smooth fibre of algebraic dimension zero there are none. In a
smooth fibre of algebraic dimension one they are exactly the translates
of its unique elliptic kernel $K_t$. In an abelian fibre they are exactly
the translates of elliptic subtori, together with the irreducible zero
divisors of the positive Appell--Humbert theta sections described below.
These descriptions, \ref{fam:src:D.1.6}, and \ref{fam:src:D.2.1} constitute a classification including
every smooth exceptional fibre and every compact irreducible surface.
\end{theorem}
\Uses{\cref{fam:src:D.1.6,fam:src:D.2.1,fam:src:D.2.2}; \cite[Theorem 4.1 and equation (4.3), pp.~1--2]{Schnell19}}
\begin{proof}
\label{app:rederive:fam-src-D.2.3}
All curves are vertical by \ref{fam:src:D.2.1}, so the list of fibre types is exhaustive.
The dimension-zero and dimension-one assertions, including uniqueness,
were proved in \ref{fam:src:D.1.6}. Let now $F_t=\mathbb C^2/\Pi\Lambda$ be abelian.
For every positive definite Hermitian form $H$ whose imaginary part takes
integer values on $\Pi\Lambda$, and every function
$\alpha:\Pi\Lambda\to U(1)$ satisfying
\[
 \alpha(\lambda+\nu)=\alpha(\lambda)\alpha(\nu)
                 e^{\pi i\operatorname{Im}H(\lambda,\nu)},
\]
take the nonzero entire solutions $\theta:\mathbb C^2\to\mathbb C$ of
\begin{equation}\label{fam:eq:cd-theta-sections}
 \theta(\zeta+\lambda)=\alpha(\lambda)
 e^{\pi H(\zeta,\lambda)+\pi H(\lambda,\lambda)/2}\theta(\zeta).
\end{equation}
The cocycle identity follows by multiplying the displayed factors; its
integral imaginary cross-term has exponential one. It defines a
holomorphic line bundle by quotienting $\mathbb C^2\times\mathbb C$.
The invariant zero set of $\theta$ descends to a compact divisor on
$F_t$. Retain precisely those zero divisors that are reduced and
irreducible. This is an exact parametrization, with all periods and
multiplier equations given; multiplying $\theta$ by a nonzero scalar
does not change its divisor.

The correspondence \cite[Theorem 4.1, p.~1, and equation (4.3), p.~2]{Schnell19} proves
that every holomorphic line bundle and every section on this actual
quotient has this form. For the cohomology matrix the exact comparison is
$H_E=\widetilde E(\,\cdot\,,J\,\cdot\,)-i\widetilde E$ from \ref{fam:src:D.1.6};
we choose the sign of $E$ making $H_E$ positive. Its imaginary part is
$-\widetilde E$, which gives the integral semicharacter equation displayed
above. A divisor class with positive Hermitian rank two gives exactly
these positive theta sections. A divisor of Hermitian rank one is pulled
back from its elliptic quotient and, when irreducible and reduced, is a
single fibre. Thus it is a translate of an elliptic subtorus. The null
lattice is computed by the explicit integer kernel in \ref{fam:src:D.1.6}. No other
rank is possible for an effective nonempty divisor, by positivity of its
averaged current. This proves exhaustiveness on the abelian fibre.

The same equations give the curves in the product covers of $S_j$:
compute them on its two-elliptic-curve lattice and take images under the
specified finite morphism. Conversely each curve downstairs is the image
of a component of its finite inverse image, so that this procedure omits
none. In $W_c$ the analogous procedure uses the homogeneous polynomial
parametrization of the normalization in \ref{fam:src:D.2.2}, with the three boundary
pairings already identified. The maps used throughout are finite or
closed embeddings, so they preserve compactness and give one-dimensional
images; no assertion is transferred from a normalization in an invalid
direction.
\end{proof}
